%======================================================================
\section{Further Remarks}\label{sec:remarks}

This section collects supplementary results. They are not needed for
the main line of the paper.

\subsection{How far can the base logic be weakened?}\label{sec:weak}
Theorem~\ref{thm:main} needs only $I\in L$ and $\St(\alpha)\in L$. In the
implicational case, the last condition is guaranteed by $\mathbf{BCI}$
(Lemma~\ref{lem:sound}). We show that none of $I$, $B$, $C$ can be
dropped. We write $\mathbf{XY}\dots$ for the implicational logic
axiomatized by the axiom schemes $X,Y,\dots$ among
\[
B=(b\to c)\to(a\to b)\to a\to c,\quad
B'=(a\to b)\to(b\to c)\to a\to c,\quad
C=(a\to b\to c)\to b\to a\to c,
\]
$K=a\to b\to a$ and $I=a\to a$, with modus ponens and substitution.

\paragraph{The axiom $I$ cannot be dropped.}
The proof uses $I$ only to discharge the premises $\gamma\to\gamma$ produced
by the decoding (Lemma~\ref{lem:key}). The next proposition shows that this
use is essential, already over $\mathbf{BC}$.

\begin{proposition}\label{prop:needI}
Let $M=(\{0,1,2\},\to,\{0\})$ with
\[
\begin{array}{c|ccc}
\to&0&1&2\\\hline
0&0&1&1\\
1&0&0&0\\
2&0&1&1
\end{array}
\]
and let $L_M$ be the set of formulas that take the value $0$ under every
valuation in $M$. Then $L_M$ is a logic containing $\mathbf{BC}$ but not $I$.
Let $\alpha=(p\to q)\to(p\to q)$. Then:
\begin{enumerate}
\item[(a)] $\alpha$ is minimal in $L_M$;
\item[(b)] $\St(\alpha)=(x_1\to(p\to q))\to((p\to q)\to x_2)\to((x_1\to x_2)\to
  x_0)\to x_0\in L_M$;
\item[(c)] $\St(\alpha)$ is not minimal in $L_M$: its one-step
  generalization
  \[
  \alpha'=(x_1\to z)\to(z\to x_2)\to((x_1\to x_2)\to x_0)\to x_0 ,
  \]
  obtained by replacing both bodies $p\to q$ by $z$, belongs to $L_M$.
\end{enumerate}
So Theorem~\ref{thm:main} fails for logics not containing $I$, even when
$\St(\alpha)\in L$.
\end{proposition}
\begin{proof}
$L_M$ is closed under substitution. It is closed under modus ponens because
$0\to y=0$ forces $y=0$. A finite check shows that $B$ and $C$ take only the
value $0$, while $2\to2=1$, so $I\notin L_M$.

(a) We have $\alpha\in L_M$: the values of implications lie in $\{0,1\}$,
and $0\to0=1\to1=0$. The one-step generalizations of $\alpha$ are
\[
(z\to q)\to(p\to q),\quad (p\to z)\to(p\to q),\quad (p\to q)\to(z\to q),\quad
(p\to q)\to(p\to z),
\]
\[
z\to(p\to q),\quad (p\to q)\to z,\quad z\to z,\quad z .
\]
Each of them takes a value $\ne0$ under some valuation in $M$. For $z\to z$
take $z=2$, and for $z$ take $z=1$. For each of the remaining six, a
refuting valuation is found by a finite check.

(b) and (c) are finite checks as well. In fact $\alpha'$ is the type of the
linear term $\lambda fgh.\,h(\lambda y.\,g(fy))$, and $\St(\alpha)$ is its
instance under $z\mapsto p\to q$. Both are theorems of $\mathbf{BC}$.
\end{proof}

Decoding $\alpha'$ as in Lemma~\ref{lem:key} gives $\beta=z\to z$. The
premise that would have to be discharged is exactly $z\to z$, which is not
available in $L_M$. Validity of $\alpha$ provides $X\to X$ only for formulas
$X$ whose value is that of an implication.
\plot{Checks: \texttt{tools/bc\_alpha\_matrix.py} (matrix $M$, found with z3
and verified by exhaustive evaluation) and \texttt{tools/bc\_check.py}
(condensed detachment: $\St(\alpha),\alpha'\in\mathbf{BC}$).}


\begin{proposition}\label{prop:weak}
\begin{enumerate}
\item[(a)] \emph{(Without $B$.)} $K$ is minimal in $\mathbf{CK}$, but
  $\St(K)=((b\to a)\to x_1)\to((a\to x_1)\to x_0)\to x_0\notin\mathbf{CK}$.
\item[(b)] \emph{(Without $C$.)} $I$ is minimal in $\mathbf{BB'KI}$, but
  $\St(I)=((a\to a)\to x)\to x\notin\mathbf{BB'KI}$. Likewise $K$ is minimal
  in $\mathbf{BK}$ and in $\mathbf{BB'KI}$, but
  $\St(K)\notin\mathbf{BB'KI}\supseteq\mathbf{BK}$.
\end{enumerate}
Hence, over $\mathbf{CK}$, $\mathbf{BK}$ and $\mathbf{BB'KI}$ (and every
logic between $\mathbf{BK}$ and $\mathbf{BB'KI}$), $\St$ does not preserve
minimality, and $L+\alpha\ne L+\St(\alpha)$ for $\alpha=K$ (resp.\ $I$).
\end{proposition}
\begin{proof}
\emph{Minimality.} The one-step generalizations of $K$ are $z\to b\to a$,
$a\to b\to z$, $a\to z$ and $z$; those of $I$ are $z\to a$, $a\to z$ and
$z$. None of them is classically valid, so none lies in any of the logics
considered, all of which are contained in $\CL$.

(a) Let $M_4=(\{0,1,2,3\},\to,\{0\})$ with the table
\[
\begin{array}{c|cccc}
\to&0&1&2&3\\\hline
0&0&1&2&3\\
1&0&0&1&0\\
2&0&0&0&3\\
3&0&1&2&0
\end{array}
\]
A finite check shows that $M_4$ validates $C$ and $K$ (hence also $I$, which
is derivable in $\mathbf{CK}$), but not $B$. The set $\{0\}$ is closed under
modus ponens, since $0\to y=0$ forces $y=0$. So every $\mathbf{CK}$-theorem
takes the value $0$. With $a=b=1$, $x_1=2$ and $x_0=3$, however,
$b\to a=0$, $(b\to a)\to x_1=2$, $a\to x_1=1$ and $(a\to x_1)\to x_0=0$.
Hence $\St(K)$ takes the value $2\to(0\to3)=2\to3=3\ne0$. Intuitively,
$\St(K)$ needs the composition $\lambda y.\,e_1(\lambda z.\,y)$ inside an
argument, which the combinators $\mathsf C,\mathsf K$ cannot form. By
contrast, $\St(I)=((a\to a)\to x)\to x$ is a $\mathbf{CK}$-theorem (it is the
type of $\mathsf C\mathsf I\mathsf I$).

(b) Let $M_3=(\{0,1,2\},\to,\{0\})$ with $x\to y=0$ if $y\le x$ and
$x\to y=1$ otherwise. $M_3$ validates $B,B',K,I$: this is a finite check.
The set $\{0\}$ is closed under modus ponens, since $0\to y=0$ forces $y=0$.
So every theorem of $\mathbf{BB'KI}$ takes the value $0$ in $M_3$. With
$a=0$ and $x=x_0=2$, however, $\St(I)$ takes the value $(0\to2)\to2=1\to2=1$.
$\St(K)$ is refuted similarly. ($M_3$ refutes $C$ and $W$.)
\end{proof}
\plot{Checks: \texttt{tools/ck\_z3.py} (matrix $M_4$, found with z3 and
verified by exhaustive evaluation), \texttt{tools/bki\_matrix.py} (matrix $M_3$),
\texttt{tools/ck\_cd.py} and \texttt{tools/bk\_cd.py} (condensed-detachment
search, consistent with (a) and (b)). Open: ticket entailment
$\mathbf{T}_\to=\mathbf{BB'IW}$, where $M_3$ does not apply because it
refutes $W$; and whether reordering the links helps. For (b), every order
of links fails already for $\St(I)$, which has a single link.}



\subsection{Additive connectives over $\mathbf{FL_e}$}\label{sec:additive}
\begin{remark}[Additive connectives over $\mathbf{FL_e}$]\label{rem:wrapper}
The proof of Theorem~\ref{thm:main} uses only (C1)--(C3), with
$I\in L$. Polarity is used only in
Lemma~\ref{lem:sound}, which is where $\mathbf{FL_e}$ with $\wedge,\vee$
fails. Two facts, both checked with a cut-free prover for $\mathbf{FL_e}$
(\texttt{tools/fle\_check.py}), show that the problem cannot be removed by
simply adding Jeřábek's ``$\wedge1$'':
\begin{enumerate}
\item[(a)] \emph{Without wrappers, $\St$ need not be sound.}
  $\alpha=(q\to q)\wedge(r\to r)$ is minimal in $\mathbf{FL_e}$, but
  $\St(\alpha)=((q\to q)\to x_0)\to((r\to r)\to x_1)\to((x_0\wedge x_1)\to
  x)\to x$ is not provable in $\mathbf{FL_e}$: each branch of the
  $\wedge$-right rule would have to discard the other branch's link.
\item[(b)] \emph{With wrappers, minimality is always lost.} Jeřábek's
  $\varphi^+$ puts $(E_u\wedge 1)$ in place of $E_u$. The link of the root is
  never discarded, so replacing its $1$ by $z$ gives a one-step
  generalization that is still provable (by $\wedge$-left). Hence
  $\varphi^+$ is not minimal in any $L\supseteq\mathbf{FL_e}$ (including
  $\mathbf{FL_{ew}}$, $\IL$, $\CL$) whenever $\varphi$ is compound. For
  example, $(((p\to p)\to x)\wedge z)\to x$ is provable.
\end{enumerate}
Exhaustive tests over all formulas with at most four binary connectives in
$\to,\wedge,\vee$ and three variables give the following. There are $413$
formulas minimal in $\mathbf{FL_e}$, up to renaming. For $164$ of them
$\St(\alpha)$ is provable, and in all $164$ cases it is minimal, as
Theorem~\ref{thm:main} predicts. $\varphi^+$ is provable in all $413$ cases
and minimal in none. Wrapping only the links strictly below an additive
occurrence gives a provable formula in all $413$ cases, but a minimal one in
only $255$. For $\alpha=p\vee(q\to q)$, for instance, the wrapper can be
removed because the $\vee$-right rule selects one branch. So a
minimality-preserving translation over $\mathbf{FL_e}$, if one exists, must
adapt the wrappers to the actual proof. Whether such a translation exists,
and the non-commutative case $\mathbf{FL}$, are left open.
\end{remark}


\subsection{Remarks on the proofs}

\begin{remark}
Corollary~\ref{cor:criterion} replaces \cite[Props.~3.1--3.2]{NM21}, whose
proofs there are sketched, and it does not need simple formulas.
\end{remark}

\begin{remark}
In Lemma~\ref{lem:sound}(i), linear $\lambda$-terms can be given instead of
sequent derivations. Without exchange the premises of $\St(\alpha)$ cannot be
reordered; see Section~\ref{sec:weak}.
\end{remark}

\subsection{Top-level premises of order 3}\label{sec:order3}

For $F=A_1\to\dots\to A_n\to r$ with $r$ a variable, we call
$A_1,\dots,A_n$ the \emph{top-level premises} of $F$. The next proposition
shows that an order-$4$ counter-example needs at least two top-level
premises of order~$3$. This explains, in part, why $\St(G')$ is long.

\begin{proposition}\label{prop:one-premise}
Let $F=A_1\to\dots\to A_n\to r\in\CL$ with $r$ a variable and
$\Ord(F)\le4$. If at most one $A_i$ has order $3$, then $F\in\IL$ or
$\IL+F=\CL$. Minimality of $F$ is not assumed.
\end{proposition}
\begin{proof}
A formula of order $\le2$ has the form $p_1\to\dots\to p_k\to q$ with
variables $p_i,q$ ($k\ge0$). We call it a \emph{Horn clause} and read it as
the rule ``from $p_1,\dots,p_k$ infer $q$''. Since $\Ord(F)\le4$, every
$A_i$ has order $\le3$. Recall that premises can be permuted in $\IL$.

Let $\Gamma$ be the set of top-level premises of order $\le 2$. For a set
$X$ of variables, let $C(X)$ be the closure of $X$ under the rules in
$\Gamma$. We identify a set of variables with the valuation making exactly
those variables true. Then:
\begin{enumerate}
\item[(H1)] $C(X)$ satisfies $\Gamma$;
\item[(H2)] if $q\in C(X)$, then $\Gamma,X\vdash_{\IL}q$ (modus ponens along
  the closure). In particular, $q\in C(\{p_1,\dots,p_k\})$ implies
  $\Gamma\vdash_\IL p_1\to\dots\to p_k\to q$.
\end{enumerate}

\emph{No premise has order $3$.} Then $F$ is $\Gamma\to r$. By (H1) and
$F\in\CL$, we get $r\in C(\emptyset)$, so $F\in\IL$ by (H2).

\emph{Exactly one premise, $A=B_1\to\dots\to B_m\to s$, has order $3$.}
Then each $B_j$ has order $\le2$. Assume $F\notin\IL$. Put
$C_0=C(\emptyset)$.
\begin{itemize}
\item $r\notin C_0$, since otherwise $\Gamma\vdash_\IL r$ by (H2).
\item Some $B_j$ is not derivable from $\Gamma$ in $\IL$. Otherwise
  $\Gamma,A\vdash_\IL s$. Moreover, the valuation $C(\{s\})$ satisfies
  $\Gamma$ and $A$, so $r\in C(\{s\})$ because $F\in\CL$. Then
  $\Gamma,s\vdash_\IL r$ by (H2), and $F\in\IL$.
\end{itemize}
Fix such a $B_j=p_1\to\dots\to p_k\to q$, and put
$C_1=C(\{p_1,\dots,p_k\})$. Then $C_0\subseteq C_1$, and $q\notin C_1$ by
(H2). The valuation $C_1$ falsifies $B_j$, hence satisfies $A$. It also
satisfies $\Gamma$ by (H1). Since $F\in\CL$, we get $r\in C_1$.

Let $H_3=\{0,\tfrac12,1\}$ be the three-element Gödel algebra, with
$x\Rightarrow y=1$ if $x\le y$ and $x\Rightarrow y=y$ otherwise. Define
$v(t)=1$ if $t\in C_0$, $v(t)=\tfrac12$ if $t\in C_1\setminus C_0$, and
$v(t)=0$ otherwise. A Horn clause $p_1\to\dots\to p_k\to q$ has value $1$
under $v$ iff $\min_i v(p_i)\le v(q)$ (the minimum of an empty set is $1$).
\begin{itemize}
\item Every clause in $\Gamma$ has value $1$, since $C_0$ and $C_1$ are
  closed under the rules in $\Gamma$.
\item $B_j$ has value $0$, since $v(p_i)\ge\tfrac12$ for all $i$ and
  $v(q)=0$. Hence $A$ has value $1$.
\item $v(r)=\tfrac12$, since $r\in C_1\setminus C_0$.
\end{itemize}
So $F$ has value $\tfrac12$, and $H_3\not\models F$. By Jankov's criterion
\cite[Theorem~1]{Jan68} (see \cite[Sect.~3.2]{NM21}), an implicational
$F\in\CL\setminus\IL$ with $H_3\not\models F$ satisfies $\IL+F=\CL$.
\end{proof}

\begin{remark}\label{rem:many}
In $\St(\alpha)$, a positive link $b_u\to x_u$ has order $3$ and a negative
link $x_u\to b_u$ has order $2$. So the top-level premises of order $3$ in
$\St(\alpha)$ are exactly the links of the positive implication occurrences
of $\alpha$. Hence, by Proposition~\ref{prop:one-premise}, if $\alpha\in\CL$
has at most one positive implication occurrence, then
$\IL+\alpha=\IL+\St(\alpha)$ is $\IL$ or $\CL$.

$\St(G')$ has $11$ positive links:
$E_1,E_3,E_5,E_7,E_9,E_{11},E_{12},E_{14},E_{15},E_{17},E_{20}$. By
Proposition~\ref{prop:one-premise}, a $\CL$-minimal axiom of order $4$ for a
proper intermediate logic has at least two top-level premises of order $3$.
The least possible number is open.
\end{remark}
